\documentclass[a4paper,12pt, openany]{book}
\usepackage[subpreambles=true]{standalone}

%----------------------------------------
% LANGUE & TYPO
%----------------------------------------
\usepackage[french]{babel}
\usepackage{csquotes}
\usepackage{fontspec}
\usepackage{unicode-math}
\usepackage{microtype}

% Paragraphe : on garde l'indentation + un petit espace entre paragraphes
\setlength{\parindent}{20pt}
\setlength{\parskip}{0.8em}

% Couleurs (si besoin)
\usepackage{xcolor}

%----------------------------------------
% POLICES
%----------------------------------------
%\setmainfont{Times New Roman}
%\setmathfont{STIX Two Math}

% Désactive toutes les small caps (évite les warnings TU/.../sc)
%\makeatletter
%\DeclareRobustCommand{\textsc}[1]{#1}
%\renewcommand{\scshape}{\upshape}
%\makeatother

%----------------------------------------
% MISE EN PAGE
%----------------------------------------
\usepackage{geometry}
\geometry{margin=2.5cm}

\usepackage{setspace}
\onehalfspacing % (si tu veux simple interligne : commente cette ligne)

\usepackage{titlesec}
\usepackage{fancyhdr}
\usepackage{ragged2e}
\justifying
\usepackage{indentfirst}

\usepackage{booktabs, array, multirow, tablefootnote}
\usepackage{caption}
\captionsetup{
  labelfont={bf},      % label en gras (pas de small caps)
  textfont={normal},
  format=hang,
  margin=12pt
}

\usepackage{float}
\usepackage{wrapfig}
\usepackage{placeins}
\usepackage{tikz}
\usepackage{tikz-3dplot}
\usetikzlibrary{positioning,arrows.meta}
\usepackage{xurl}
\usepackage{epigraph}
\usepackage{enumitem}
\usepackage{comment}
\usepackage{mwe} % à retirer si inutile

%----------------------------------------
% GRAPHISMES & LIENS
%----------------------------------------
\usepackage{graphicx}

%----------------------------------------
% BIBLIO
%----------------------------------------
\usepackage[
  backend=biber,
  style=apa
]{biblatex}
\addbibresource{../../ref.bib}

% Mise en forme biblio
\renewcommand*{\bibfont}{\normalsize}
\setlength\bibhang{1cm}
\setlength\bibitemsep{6pt}
\defbibheading{bibliography}{\raggedright \textbf{Références} \vspace{12pt}}

% Neutralise l'usage potentiel des small caps par biblatex (selon styles)
\AtBeginDocument{%
  \providecommand*{\mkbibnamefamily}[1]{#1}%
  \renewcommand*{\mkbibnamefamily}[1]{#1}% noms d'auteurs sans \textsc
}

%----------------------------------------
% DATE
%----------------------------------------
\usepackage[french]{datetime2}

%----------------------------------------
% HYPERREF EN DERNIER
%----------------------------------------
\usepackage[hidelinks]{hyperref}

%----------------------------------------
% PIEDS DE PAGE
%----------------------------------------
\pagestyle{fancy}
\fancyhf{}
\fancyfoot[C]{\thepage}
\renewcommand{\headrulewidth}{0pt}
\renewcommand{\footrulewidth}{0pt}

% Veuves/orphelines
\clubpenalty=100000
\widowpenalty=100000
\displaywidowpenalty=100000
\brokenpenalty=100000



\begin{document}

%\section{La relation usager : un cas particulier}

\paragraph{Arrivée du NPM et ce que c'est} A

À la fin des années 1970, à la suite d’une récession économique, apparaît un nouveau courant de pratiques de gestion de l’administration : le \textit{New Public Management} (NPM). Ce courant émerge au Royaume-Uni sous Margaret Thatcher, puis localement aux États-Unis, avant de s’étendre à la Nouvelle-Zélande et à l’Australie \parencite{gruening01}. Le NPM emprunte des méthodes et des approches au secteur privé, inspirées des théories économiques des marchés, et les applique au secteur public \parencite{pollitt11}. L’un des principes centraux du NPM est de considérer les citoyens comme des clients, notamment en privilégiant ce qu’ils souhaitent \parencite{aberbach05}. La publication de \textit{Reinventing Government: How the Entrepreneurial Spirit Is Transforming the Public Sector} par \textcite{osborne92} \parencite{thomas13} ainsi, que les premiers succès des politiques inspirées du marché \parencite{bashevkin94} encouragent de nombreux gouvernements à adopter cette perspective, renforçant encore l’influence du NPM \parencite{haque98, kaul96}. \par


\paragraph{Cependant la relation Etat-citoyen repose sur un contrat social} B

Si le NPM tend à assimiler les citoyens à des clients, la relation qui unit l’individu à l’État s’inscrit avant tout dans un \textit{contrat social} fondé sur une obligation réciproque : l'État doit des services et droits au citoyen en échange d'obéissance et de devoirs \parencite{haque99b}. Cette relation requiert une transparence de l'État sur ses processus plus exigeante que lors d'une transaction mercantile, souvent évaluée seulement sur le résultat. \textcite{behn98} explique la contradiction entre l'exigence de résultats et l'exigence de processus ; l'exigence de résultats demande des objectifs clairs et précis. Or, les décideurs publics répondent à une grande diversité d'audiences avec des intérêts différents, si ce n'est mutuellement exclusifs, ils sont donc fortement incités à rester ambigu. Ainsi, la nature même de la relation entre l'administration et les citoyens et la contradiction entre la transparence de processus et celle de résultats sont une limite au modèle «\;usager = client\;». À cela s'ajoutent les spécificités des responsabilités de l'État. \par


\paragraph{Responsabilités de processus de l'État : incompatibles avec NPM} C

Les responsabilités de l’État qui diffèrent des logiques managériales portées par le NPM peuvent être déclinées en deux catégories : les responsabilités de procédure et les responsabilités de nature de service. Alors que le NPM met l’accent sur la performance, la compétition et la satisfaction individuelle, les missions de l’État impliquent la garantie de l’égalité d’accès, la continuité des services publics, la justice sociale et la protection des droits fondamentaux \parencite{hood91, haque99}. Ces objectifs nécessitent de trouver un équilibre entre traiter les demandes équitablement et impartialement tout en les traitant au cas par cas lors de situations plus complexes \parencite{brewer07}. Cet équilibre est difficilement conciliables avec une logique clientéliste fondée sur la segmentation, le choix et les mécanismes de marché. Les travaux comparatifs de \textcite{pollitt11} montrent que la transposition des méthodes du secteur privé au secteur public engendre des tensions entre performance mesurable et valeurs publiques non quantifiables, comme l’équité ou la transparence. 

\paragraph{Responsabilités de nature de service de l'État : incompatibles avec NPM} D

Au-delà des processus utilisés pour l'offre de service et d'une obligation de transparence, la relation entre l'État et le citoyen repose également sur la nature des services offerts. L'obligation de l'État de fournir des services basiques au citoyen est reconnue dans la plupart des sociétés \parencite{haque99}. Or, l'État possède une emprise sur plusieurs de ces services incompatible avec une relation marchande traditionnelle. Cette emprise se manifeste de deux manières différentes :
\begin{itemize}
	\item \textbf{Les services sont non rivaux} ce qui donne un monopole légal, l'usager n'a pas le choix de son fournisseur de service \parencite{christensen02}, par exemple il est illégal de dispenser de la justice à la place de l'État.
	\item \textbf{Les services sont non exclusifs} ce qui empêche les usagers de refuser un service comme il est illégal de refuser de payer ses impôts.
\end{itemize}
Cette emprise de l'État empêche ses usagers de recourir à l'option de se désengager (Exit). Pourtant sans cette option il ne peut y avoir de rivalité ce qui se situe au cœur de toute logique de marché.


\paragraph{Impossibilité de sortie renforce l'importance de la réclamation usager} E










\textcite{brewer07} Les citoyens se plaignent en raison d'expériences inappropriées, désobligeantes, d'omissions, d'erreurs, d'incohérences, de conseils et procédures prêtant à confusion ou de biais et d'injustices venant d'agents publics.

Réclamations comme outil de responsabilisation 













%\printbibliography



\end{document}